
\pagestyle{empty}  
\section*{Declaración de originalidad}

Declaro bajo mi responsabilidad que el Proyecto presentado con el título
\textbf{"Separación de señales electromiográficas según el músculo de procedencia"} en la ETS de Ingeniería – ICAI de la Universidad Pontificia
Comillas en el curso académico 2025-2026 es de mi autoría y no ha
sido presentado con anterioridad a otros efectos.
El Proyecto no es plagio de otro, ni total ni parcialmente, y la
información que ha sido tomada de otros documentos está debidamente
referenciada.

\vspace{3cm}

\noindent
\begin{tabular}{@{}p{0.45\textwidth}@{}}
  Firmado (alumno): \\[2.5cm]
  \hrule \\[0.3cm]
  Antón Cobian Iregui \\[1cm]
  Fecha: \fechaDeclaracion
\end{tabular}

\newpage



\section*{Uso de Inteligencia Artificial\footnote{%
  Esta declaración se refiere al uso de la Inteligencia Artificial
  generativa para realizar los documentos del Proyecto (Anexo~B y
  Memoria). No aplica a Proyectos donde, por su naturaleza, deban
  emplear inteligencia artificial como parte de los mismos (aplicación
  de técnicas de aprendizaje automático, redes neuronales, análisis de
  datos…).}}

Declaro bajo mi responsabilidad que (\emph{indicar la opción correcta}):

\bigskip

\begin{itemize}[leftmargin=2em, label=$\square$]

  \item No he utilizado Inteligencia Artificial en la elaboración del
        presente documento.

  \bigskip

  \item He utilizado Inteligencia Artificial en la elaboración del
        presente documento y/o del Anexo~B siempre en las condiciones
        permitidas por la Universidad Pontificia Comillas, es decir,
        aplicando el Nivel~2 de la
        \href{https://aiassessmentscale.com/}{Escala de Evaluación de
        Perkins et al.\ (2024)}: \emph{«La IA puede utilizarse para
        actividades previas a la tarea, como la lluvia de ideas, la
        descripción y la investigación inicial. Este nivel se centra en
        el uso de la IA para la planificación, las síntesis y la
        generación de ideas, pero las evaluaciones deben hacer hincapié
        en la capacidad de desarrollar y refinar estas ideas de forma
        independiente».}

        En concreto, la Inteligencia Artificial ha sido empleada para:

        \bigskip
        \noindent
        \begin{tabularx}{\textwidth}{|X|}
          \hline
          \rule{0pt}{2.5cm}%
          \textit{(indicar aquí el uso concreto que se ha hecho de la
          Inteligencia Artificial)} \\
          \hline
        \end{tabularx}

\end{itemize}

\vspace{2cm}

\noindent
\begin{tabular}{@{}p{0.45\textwidth}@{}}
  Firmado (alumno): \\[2.5cm]
  \hrule \\[0.3cm]
  Antón Cobian Iregui \\[1cm]
  Fecha: \fechaDeclaracion
\end{tabular}

\newpage


% ============================================================
%  3. AUTORIZACIÓN PARA LA ENTREGA DEL PROYECTO
% ============================================================
\section*{Autorización para la entrega del Proyecto}

\bigskip

\noindent
\begin{tabularx}{\textwidth}{|X|X|}
  \hline
  \textbf{El Director del Proyecto}
    & \textbf{El Co-Director del Proyecto} (\textit{si aplica}) \\[4cm]
  Fdo: \hrulefill & Fdo: \hrulefill \\[0.5cm]
  Fecha: \hrulefill & Fecha: \hrulefill \\[0.5cm]
  \hline
\end{tabularx}

